\chapter{Reproducible Commands and Runbook}
\label{app:commands}

This appendix lists the commands needed to reproduce the thesis runs. The defended pipeline, tests, experiment tracking, and figure generation run in the \texttt{thesis-mlops} conda environment. A separate \texttt{tf-gpu} environment was used only to accelerate model building; it is not required for normal pipeline inference.

\section*{A.1 Environment setup}
\begin{footnotesize}
\begin{verbatim}
conda create -n thesis-mlops python=3.11 -y
conda activate thesis-mlops
python -m pip install --upgrade pip
python -m pip install -r config/requirements.txt
\end{verbatim}
\end{footnotesize}

\section*{A.2 Retrieve versioned data and model artifacts}
\begin{footnotesize}
\begin{verbatim}
dvc pull
\end{verbatim}
\end{footnotesize}

\section*{A.3 Run the assessment pipeline (Stages 1--7)}
For a specific raw session, pass the accelerometer and gyroscope CSV files explicitly:
\begin{footnotesize}
\begin{verbatim}
python run_pipeline.py --stages ingestion validation transformation \
  inference evaluation monitoring trigger \
  --accel-file data/raw/<session>_accelerometer.csv \
  --gyro-file  data/raw/<session>_gyroscope.csv
\end{verbatim}
\end{footnotesize}

\section*{A.4 Governed retraining (Stages 8--10, review-gated)}
\begin{footnotesize}
\begin{verbatim}
python run_pipeline.py --retrain --labels data/raw/labels.csv
\end{verbatim}
\end{footnotesize}
Retraining fine-tunes the current proposed model on labeled data and requires review, labels, and quality-assurance approval before any promotion; it is never triggered automatically.

\section*{A.5 Chapter~7 figures}
\begin{footnotesize}
\begin{verbatim}
python scripts/generate_ch7_figures.py
\end{verbatim}
\end{footnotesize}

\section*{A.6 Tests}
\begin{footnotesize}
\begin{verbatim}
python -m pytest tests -q
python -m pytest tests/test_runtime_contract_consistency.py -q
python -m pytest tests/test_stage2_validation.py tests/test_stage6_monitoring.py -q
\end{verbatim}
\end{footnotesize}

\section*{A.7 API inference service (local)}
\begin{footnotesize}
\begin{verbatim}
python -m uvicorn src.api.app:app --host 127.0.0.1 --port 8000
\end{verbatim}
\end{footnotesize}
